\documentclass[11pt,a4paper]{article}
\usepackage[utf8]{inputenc}
\usepackage{amsmath,amssymb,amsthm}
\usepackage{mathrsfs}
\usepackage{geometry}
\geometry{margin=1in}
\usepackage{hyperref}
\usepackage{enumitem}

\newtheorem{definition}{Definition}[section]
\newtheorem{theorem}{Theorem}[section]
\newtheorem{proposition}{Proposition}[section]
\newtheorem{remark}{Remark}[section]
\newtheorem{corollary}{Corollary}[section]

\title{\textbf{Spectral Theory and Quantum Symmetries\\in the Algebraic Framework}\\[6pt]
\large Lecture by Valter Moretti}
\author{Lecture Notes from the AQFT 50th Anniversary Conference}
\date{}

\begin{document}
\maketitle

\begin{abstract}
Valter Moretti presents the spectral-theoretic and group-theoretic foundations
relevant to algebraic quantum field theory. He discusses the rigorous formulation
of quantum symmetries via Wigner's theorem, the role of the spectral theorem for
unbounded operators in the algebraic framework, the construction of the one-particle
space and its relation to the Fock space, and the interplay between modular theory
and the KMS condition. Applications to quantum field theory on curved spacetimes,
including the definition of Hadamard states through spectral methods, are reviewed.
\end{abstract}

\tableofcontents

%---------------------------------------------------------------------
\section{Quantum Symmetries: Wigner's Theorem}
%---------------------------------------------------------------------

\subsection{States and Symmetries}

In quantum mechanics, a \emph{state} is a ray $[\psi] = \{e^{i\alpha}\psi \mid
\alpha \in \mathbb{R}\}$ in a complex Hilbert space $\mathcal{H}$. The transition
probability between two states is:
\begin{equation}
  P([\psi], [\phi]) = \frac{|\langle \psi, \phi \rangle|^2}
  {\|\psi\|^2 \|\phi\|^2}.
\end{equation}

A \emph{symmetry} is a bijection $s$ on the set of rays preserving
transition probabilities: $P(s[\psi], s[\phi]) = P([\psi], [\phi])$.

\begin{theorem}[Wigner, 1931]
Every symmetry $s$ is implemented by either a \textbf{unitary} operator
$U$ or an \textbf{anti-unitary} operator $A$:
\[
  s[\psi] = [U\psi] \quad\text{or}\quad s[\psi] = [A\psi].
\]
The implementing operator is unique up to a phase factor $e^{i\alpha}$.
\end{theorem}

\subsection{Kadison's Generalization}

Kadison extended Wigner's theorem to the algebraic setting: a symmetry
is an affine bijection on the convex set of density matrices (mixed states).
Every such map is implemented by a Jordan $*$-automorphism, which is either
a $*$-automorphism or a $*$-anti-automorphism.

\subsection{Continuous Symmetry Groups}

For a continuous symmetry group $G$ (e.g., the Poincar\'e group), a
\emph{projective unitary representation} arises:
\begin{equation}
  U(g_1)\, U(g_2) = e^{i\omega(g_1, g_2)}\, U(g_1 g_2),
\end{equation}
where $\omega: G \times G \to \mathbb{R}$ is a 2-cocycle. Bargmann's theorem
(1954) states that for simply connected Lie groups with $H^2(\mathfrak{g}, \mathbb{R}) = 0$
(e.g., the Poincar\'e group's universal cover), projective representations
lift to true representations of the universal covering group.

%---------------------------------------------------------------------
\section{Spectral Theory for Unbounded Operators}
%---------------------------------------------------------------------

\subsection{Self-Adjoint Operators and the Spectral Theorem}

Observables in quantum mechanics are represented by (generally unbounded)
self-adjoint operators. The spectral theorem provides the connection to
measurement:

\begin{theorem}[Spectral Theorem]
For every self-adjoint operator $A$ on $\mathcal{H}$, there exists a unique
projection-valued measure $E^A$ on $\mathbb{R}$ such that:
\begin{equation}
  A = \int_{\mathbb{R}} \lambda\; dE^A(\lambda),
  \qquad \mathcal{D}(A) = \Bigl\{\psi \in \mathcal{H} \;\Big|\;
  \int_{\mathbb{R}} \lambda^2\, d\|E^A(\lambda)\psi\|^2 < \infty\Bigr\}.
\end{equation}
\end{theorem}

\subsection{Functional Calculus}

For any Borel function $f: \mathbb{R} \to \mathbb{C}$:
\begin{equation}
  f(A) = \int_{\mathbb{R}} f(\lambda)\; dE^A(\lambda).
\end{equation}

Key applications:
\begin{itemize}[nosep]
  \item \textbf{Stone's theorem}: one-parameter unitary groups $U(t) = e^{itA}$
        are in bijection with self-adjoint operators $A$.
  \item \textbf{Modular operator}: $\Delta^{it}$ as a one-parameter group
        generated by $\log\Delta$.
  \item \textbf{KMS condition}: $\omega$ is KMS at inverse temperature $\beta$
        for the group $\alpha_t = \mathrm{Ad}(e^{itH})$ if the function
        $t \mapsto \omega(A\,\alpha_t(B))$ extends analytically to the
        strip $0 < \mathrm{Im}(t) < \beta$.
\end{itemize}

%---------------------------------------------------------------------
\section{One-Particle Structures and Fock Space}
%---------------------------------------------------------------------

\subsection{The One-Particle Hilbert Space}

For a free Klein--Gordon field on a globally hyperbolic spacetime $(M, g)$,
the one-particle Hilbert space is constructed as follows:
\begin{enumerate}[nosep]
  \item Start with the space of real solutions $\mathrm{Sol}(M)$ to the
        Klein--Gordon equation.
  \item Equip it with the symplectic form:
        \begin{equation}
          \sigma(\phi_1, \phi_2) = \int_\Sigma \bigl(\phi_1\, \nabla_n \phi_2
          - \phi_2\, \nabla_n \phi_1\bigr)\, d\Sigma,
        \end{equation}
        where $\Sigma$ is any Cauchy surface.
  \item Choose a \textbf{complex structure} $J: \mathrm{Sol}(M) \to \mathrm{Sol}(M)$
        compatible with $\sigma$ ($J^2 = -\mathrm{id}$ and $\sigma(\cdot, J\cdot) > 0$).
  \item Complete: $\mathcal{H}_1 = \overline{(\mathrm{Sol}(M), \langle\cdot,\cdot\rangle_J)}$,
        where $\langle\phi_1, \phi_2\rangle_J = \sigma(\phi_1, J\phi_2) + i\sigma(\phi_1, \phi_2)$.
\end{enumerate}

Different choices of $J$ give different (generally inequivalent) one-particle
spaces and hence different Fock representations.

\subsection{The Bosonic Fock Space}

The symmetric (bosonic) Fock space over $\mathcal{H}_1$:
\begin{equation}
  \mathcal{F}_s(\mathcal{H}_1) = \bigoplus_{n=0}^{\infty}
  \mathcal{H}_1^{\otimes_s n},
\end{equation}
with vacuum vector $\Omega = (1, 0, 0, \ldots)$, creation and annihilation
operators $a^*(f), a(f)$, and Weyl operators
$W(f) = e^{i\overline{\phi(f)}}$.

\subsection{Hadamard Condition and Complex Structures}

On curved spacetimes, the Hadamard condition selects physically meaningful
complex structures:

\begin{theorem}
A quasi-free state $\omega_J$ (determined by a complex structure $J$) is
Hadamard if and only if $J$ differs from any reference Hadamard complex
structure $J_0$ by a smoothing operator:
\[
  J - J_0 \in \Psi^{-\infty}(M).
\]
\end{theorem}

This ensures that the short-distance singularity structure of the two-point
function is universal (determined by the geometry), while the smooth part
is state-dependent.

%---------------------------------------------------------------------
\section{Modular Theory and KMS States}
%---------------------------------------------------------------------

\subsection{The Tomita--Takesaki Theorem}

For a von Neumann algebra $\mathfrak{M}$ with a cyclic and separating vector
$\Omega$, the modular operator $\Delta$ and modular conjugation $J$ satisfy:
\begin{align}
  \Delta^{it}\, \mathfrak{M}\, \Delta^{-it} &= \mathfrak{M}
  \quad\forall\, t \in \mathbb{R}, \\
  J\, \mathfrak{M}\, J &= \mathfrak{M}'.
\end{align}

\subsection{The KMS Condition}

\begin{definition}
A state $\omega$ on a $C^*$-algebra $\mathfrak{A}$ with dynamics
$\alpha: \mathbb{R} \to \mathrm{Aut}(\mathfrak{A})$ satisfies the
\textbf{KMS condition at inverse temperature $\beta > 0$} if for all
$A, B \in \mathfrak{A}$, the function
\[
  F_{A,B}(t) = \omega(A\, \alpha_t(B))
\]
extends to a function analytic in the strip $\{z \mid 0 < \mathrm{Im}(z) < \beta\}$,
continuous on its closure, with boundary condition:
\[
  F_{A,B}(t + i\beta) = \omega(\alpha_t(B)\, A).
\]
\end{definition}

\begin{theorem}[Takesaki]
The vector state $\omega_\Omega(A) = \langle\Omega, A\Omega\rangle$ is a
KMS state at $\beta = -1$ for the modular automorphism group
$\sigma_t = \mathrm{Ad}(\Delta^{it})$.
\end{theorem}

\subsection{Physical Applications}

\begin{itemize}[nosep]
  \item \textbf{Bisognano--Wichmann}: the vacuum restricted to a wedge
        algebra is KMS at $\beta = 2\pi$ for the Lorentz boost.
  \item \textbf{Unruh effect}: an accelerated observer experiences the vacuum
        as a thermal state at temperature $T = a/(2\pi)$.
  \item \textbf{Hawking radiation}: the Hartle--Hawking state on a Schwarzschild
        black hole is KMS at the Hawking temperature $T_H = \kappa/(2\pi)$.
\end{itemize}

%---------------------------------------------------------------------
\section{Algebraic Formulation of QFT}
%---------------------------------------------------------------------

\subsection{The Weyl Algebra}

For the free scalar field, the Weyl algebra $\mathcal{W}(\mathrm{Sol}, \sigma)$
is the $C^*$-algebra generated by unitaries $W(f)$, $f \in \mathrm{Sol}$,
satisfying:
\begin{align}
  W(f)^* &= W(-f), \\
  W(f)\, W(g) &= e^{-\frac{i}{2}\sigma(f,g)}\, W(f+g).
\end{align}

This algebra is independent of any choice of state or representation
and encodes the canonical commutation relations in exponentiated form.

\subsection{States Determine Physics}

Two states $\omega_1, \omega_2$ on $\mathcal{W}$ may yield unitarily
inequivalent GNS representations. The physical content (particle spectrum,
thermal properties, symmetry breaking) depends on the state.

In Minkowski spacetime, the vacuum state is selected by Poincar\'e invariance
and the spectrum condition. On generic curved spacetimes, the Hadamard
condition replaces these criteria, selecting a large but physically
meaningful class of states.

%---------------------------------------------------------------------
\section{Quantum Fields on Curved Spacetimes}
%---------------------------------------------------------------------

\subsection{The Locally Covariant Framework}

The Weyl algebra construction is functorial: for each globally hyperbolic
spacetime $(M, g)$, one obtains $\mathcal{W}(M, g)$, and isometric
causally convex embeddings $\chi: M \hookrightarrow N$ induce injective
$*$-homomorphisms $\alpha_\chi: \mathcal{W}(M) \to \mathcal{W}(N)$.

\subsection{The Stress-Energy Tensor}

The renormalized stress-energy tensor is obtained through:
\begin{enumerate}[nosep]
  \item Point-splitting regularization using the Hadamard parametrix.
  \item Subtraction of the singular (state-independent) part.
  \item Addition of geometric counterterms for conservation.
\end{enumerate}

The ambiguity is parameterized by finitely many real numbers (Wald's axioms),
corresponding to local curvature scalars of the correct dimension.

%---------------------------------------------------------------------
\section{Conclusions}
%---------------------------------------------------------------------

\begin{itemize}[nosep]
  \item Wigner's theorem and its generalizations provide the mathematical
        foundation for quantum symmetries in the algebraic framework.
  \item The spectral theorem underpins both the physical interpretation
        of observables and the mathematical structure of modular theory.
  \item One-particle structures and Fock spaces depend on a choice of
        complex structure; the Hadamard condition constrains this choice
        on curved spacetimes.
  \item The KMS condition connects modular theory to thermodynamics,
        providing a unified understanding of the Unruh and Hawking effects.
  \item The algebraic framework, being representation-independent, is
        the natural setting for quantum field theory on generic spacetimes.
\end{itemize}

\end{document}
